% Lösungen Einheit 3 von 4 – Von der Sekante zur Tangente (zusammenbau v0.1)
\einheitenkopf{Einheit 3 von 4 · Von der Sekante zur Tangente}

\erg{16}{a) Sekantensteigung – mittlere Änderungsrate: Differenzenquotient über das Intervall von $2$ bis $5$ \quad b) $m = 5$; $4{,}5$; $4{,}1$; $4{,}01$ – die Steigungen nähern sich $4$: $f'(2) = 4$ \quad c) Mittlere Steigung $m = \frac{f(2) - f(0)}{2} = \frac{-2 - 2}{2} = -2$; Tangente in $W(1 | 0)$: Steigungsdreieck 1 nach rechts, 3 nach unten, Steigung $m_W = -3$ (Vorzeichen negativ, der Graph fällt) \quad d) mittlere Rate $\frac{T(40) - T(0)}{40} \approx -0{,}249$ °C je min, momentane Rate $T'(40) \approx -0{,}225$ °C je min; Abweichung $\frac{0{,}249 - 0{,}225}{0{,}225} \approx 0{,}107 = 10{,}7\,\%$ – ja, mehr als $10\,\%$ (auf die mittlere Rate bezogen wären es nur $9{,}7\,\%$ – falsche Bezugsgröße, falsches Urteil) \quad e) Sekantensteigung $\frac{18 - 0}{3} = 6$; $f'(x_0) = 3x_0^2 - 3 = 6$ ⇔ $x_0^2 = 3$ ⇔ $x_0 = \sqrt{3} \approx 1{,}73$ (auch $-\sqrt{3}$); an dieser Stelle ist die Tangente parallel zur Sekante durch die Graphenpunkte bei $0$ und $3$ \quad f) Die Sekante durch die Schnittpunkte ist $g$ selbst, die mittlere Änderungsrate also $m = 0{,}5$ (keine Rechnung nötig); $f'(x) = \frac{1}{2\sqrt{x}} = 0{,}5$ ⇔ $\sqrt{x} = 1$ ⇔ $x = 1$ \quad g) „Bestimme alle Stellen, an denen die Tangente an den Graphen von $f$ parallel zur Sekante durch die Graphenpunkte bei $-2$ und $2$ ist.“ – links steht die Sekantensteigung, rechts die Tangentensteigung; die Lösungen sind die Stellen $x \approx -1{,}15$ und $x \approx 1{,}15$ (nicht: „Löse die Gleichung“) \quad h) Randsekante durch $(1 | 0)$ und $P$: $m = 0{,}5$; Tangentensteigung $= 0{,}25$ kommt nicht vor; links von $P$ fallen die Sekantensteigungen von $0{,}5$ gegen $0{,}25$, rechts von $0{,}25$ gegen $0$ (rechtsgekrümmt, immer flacher) – alle $m$ mit $0 < m \le 0{,}5$ außer $0{,}25$}

16c)

\begin{ksys}[xmin=-1,xmax=3,ymin=-3,ymax=4,klein]\funktion{\x^3-3*\x^2+2}{f}\tangentean{\x^3-3*\x^2+2}{1}{t}\end{ksys}

\erg{17}{a) I: Sekante durch die Graphenpunkte bei $t = 2$ und $t = 8$ – mittlere Änderungsrate des Wasserstands zwischen der 2. und der 8. Stunde in cm je Stunde (etwa $3$); II: Tangente im Graphenpunkt bei $t = 3$ – momentane Änderungsrate des Wasserstands nach 3 Stunden in cm je Stunde (etwa $5$)}
\erg{18}{a) Fehler: Die Tangentensteigung $0{,}25$ ist mitgezählt, aber die Tangente hat mit dem Graphen nur $P$ gemeinsam. Richtig: alle $m$ mit $0 < m \le 0{,}5$ außer $m = 0{,}25$. \quad b) Je näher $Q$ an $P$ liegt, desto weniger weicht der Graph zwischen $P$ und $Q$ von seiner Tangente in $P$ ab; die Sekante dreht sich in die Tangente hinein, ihre Steigung strebt gegen die Tangentensteigung – erreicht wird sie nicht, weil $Q$ nie mit $P$ zusammenfällt. \quad c) mittlere Laderate $\frac{L(10) - L(0)}{10} \approx 5{,}06$ Prozent je min, momentane Laderate $L'(10) = 8 \cdot \mathrm{e}^{-1} \approx 2{,}94$ Prozent je min; $\frac{5{,}06 - 2{,}94}{5{,}06} \approx 0{,}42$ – die momentane Rate liegt etwa $42\,\%$ unter der mittleren, das Laden wird langsamer \quad d) Sekante durch die Graphenpunkte bei $0$ und $4$ mit Steigung $1$; Tangente im Punkt $(2 | 2)$ mit Steigung $f'(2) = 1$ – beide Geraden sind parallel}

18d)

\begin{ksys}[xmin=-1,xmax=5,ymin=-1,ymax=7,klein]\parabel{0.25}{0}{1}{f}\gerade{1}{1}{s}\tangentean*{0.25*\x^2+1}{2}{t}\end{ksys}

